\labsection{5}{Stack ADT and expression processing}{sec:lab05}

\begin{sourcealignment}{Official Lab Manual --- Lab V}
\textbf{Objective.} Build intuition for stacks, their operations, implementations, and use cases.

\textbf{Outcomes.} Understand stack behavior; use push/pop/peek; apply a stack to unseen programming problems.

\textbf{Problem directions.} Implement a stack with an array; implement a stack with a linked list; convert infix notation to postfix using a stack.
\end{sourcealignment}

\begin{labproject}{Implement the stack ADT twice}
\textbf{Coverage.} Chapter 5.

Required operations are \texttt{push()}, \texttt{pop()}, \texttt{peek()}, \texttt{isFull()}, \texttt{isEmpty()}, and \texttt{display()} for the array version and the corresponding linked-list behavior.

\begin{lstlisting}
// Suggested trace
push(10);
push(20);
push(30);
peek();       // 30
pop();        // removes 30
display();    // 20 10 from top to bottom
\end{lstlisting}

\begin{tryit}
Deliberately trigger array overflow and stack underflow. Then implement the manual's infix-to-postfix exercise and test \texttt{A+B*C} and \texttt{(A+B)*C}.
\end{tryit}

\begin{clarification}
A linked-list stack normally has no predetermined fixed capacity. If your interface contains \texttt{isFull()}, document whether you impose a maximum node count; otherwise only allocation failure makes the structure unable to grow.
\end{clarification}
\end{labproject}

\begin{verification}
\begin{itemize}
  \item LIFO order is visible in the output.
  \item Underflow is handled before reading/removing the top.
  \item Array overflow is handled before writing past capacity.
  \item The infix-to-postfix result respects parentheses and operator precedence.
\end{itemize}
\end{verification}

\begin{labdeliverable}
Submit both stack implementations, one overflow/underflow trace, and the infix-to-postfix converter with at least two test expressions.
\end{labdeliverable}
